\clearpage\phantomsection\addcontentsline{toc}{chapter}{REFERENCES}\markboth{References}{}\begin{singlespace}\begin{thebibliography}{99}
\interlinepenalty=10000

\bibitem{ref1} P. Y. Wen et al., “RANO 2.0: Update to the Response Assessment in Neuro-Oncology criteria for high- and low-grade gliomas in adults,” Journal of Clinical Oncology, vol. 41, no. 33, pp. 5187–5199, 2023, doi: 10.1200/JCO.23.01059.

\bibitem{ref2} M. A. Qazi, M. Nadeem, and M. Yaqub, “Beyond generative AI: World models for clinical prediction, counterfactuals, and planning,” arXiv preprint arXiv:2511.16333, 2025, doi: 10.48550/arXiv.2511.16333.

\bibitem{ref3} M. Assran et al., “Self-supervised learning from images with a joint-embedding predictive architecture,” in Proc. IEEE/CVF Conference on Computer Vision and Pattern Recognition, 2023, pp. 15619–15629, doi: 10.1109/CVPR52729.2023.01499.

\bibitem{ref4} Y. Suter et al., “The LUMIERE dataset: Longitudinal Glioblastoma MRI with expert RANO evaluation,” Scientific Data, vol. 9, art. 768, 2022, doi: 10.1038/s41597-022-01881-7.

\bibitem{ref5} Y. Yang et al., “Medical world model,” in Proc. IEEE/CVF International Conference on Computer Vision, 2025, pp. 8319–8329, doi: 10.1109/ICCV51701.2025.00779.

\bibitem{ref6} V. Biller et al., “TumorFlow: Physics-guided longitudinal MRI synthesis of glioblastoma growth,” arXiv preprint arXiv:2603.04058, 2026, doi: 10.48550/arXiv.2603.04058.

\bibitem{ref7} Q. Liu et al., “Treatment-aware diffusion probabilistic model for longitudinal MRI generation and diffuse glioma growth prediction,” IEEE Transactions on Medical Imaging, vol. 44, no. 6, pp. 2449–2462, 2025, doi: 10.1109/TMI.2025.3533038.

\bibitem{ref8} D. Laslo et al., “Generative AI for spatial tumor growth on MRI: A proof-of-principle study in pediatric diffuse midline glioma,” BMC Medicine, vol. 24, art. 389, 2026, doi: 10.1186/s12916-026-04911-y.

\bibitem{ref9} S. Huisman, N. Belkacemi, V. C. Keil, J. Verhoeff, and S. David, “Rectified flow-based prediction of post-treatment brain MRI from pre-radiotherapy priors for patients with glioma,” arXiv preprint \href{https://arxiv.org/abs/2603.08385v2}{arXiv:2603.08385v2}, 2026, doi: 10.48550/arXiv.2603.08385.

\bibitem{ref10} Y. Yue, Y. Wang, C. Tao, P. Liu, S. Song, and G. Huang, “CheXWorld: Exploring image world modeling for radiograph representation learning,” in Proc. IEEE/CVF Conference on Computer Vision and Pattern Recognition, 2025, pp. 20778–20788.

\bibitem{ref11} H. Huang et al., “Learning sparse latent predictive foundation model for multimodal neuroimaging,” arXiv preprint \href{https://arxiv.org/abs/2606.14957v5}{arXiv:2606.14957v5}, 2026, doi: 10.48550/arXiv.2606.14957.

\bibitem{ref12} F. Mager and L. K. Hansen, “Contrastive joint-embedding prediction for representation learning in structural MRI,” arXiv preprint arXiv:2607.11962, 2026, doi: 10.48550/arXiv.2607.11962.

\bibitem{ref13} T. Z. Li et al., “Self-supervised learning of imaging and clinical signatures using a multimodal joint-embedding predictive architecture,” arXiv preprint arXiv:2509.15470, 2025, doi: 10.48550/arXiv.2509.15470.

\bibitem{ref14} J. Bai and J. Xiong, “Temporal-distance JEPA: Plan-aware representation learning for latent world model predictive control,” arXiv preprint arXiv:2607.25337, 2026, doi: 10.48550/arXiv.2607.25337.

\bibitem{ref15} L. Chen, N. Zebardast, M. Wang, T. Elze, and J. I. Comander, “Training-inference input alignment outweighs framework choice in longitudinal retinal image prediction,” arXiv preprint \href{https://arxiv.org/abs/2604.16955v1}{arXiv:2604.16955v1}, 2026, doi: 10.48550/arXiv.2604.16955.

\bibitem{ref16} D. Tak et al., “A generalizable foundation model for analysis of human brain MRI,” Nature Neuroscience, vol. 29, pp. 945–956, 2026, doi: 10.1038/s41593-026-02202-6.

\bibitem{ref17} Y. Suter et al., “LUMIERE dataset: MRI data and automated segmentations,” Figshare, version 1, 2022, data set, doi: 10.6084/m9.figshare.21249516.v1.

\bibitem{ref18} I. T. Hovden et al., “Brain tumor MRI: Serial assessments in longitudinal oncological research (SAILOR),” EBRAINS, version 1, 2023, data set, doi: 10.25493/D948-J3P.

\bibitem{ref19} D. Yaseen et al., “University of Missouri Post-operative Glioma Dataset (MU-Glioma-Post),” The Cancer Imaging Archive, version 1, 2025, data set, doi: 10.7937/7K9K-3C83.
\bibitem{record} Project research record, ``Experiments and evidence'' and ``Methods,'' evidence through November 9, 2026; ``Final study synthesis,'' closed November 11, 2026. Repository records: \texttt{info/EXPERIMENTS.md}, \texttt{info/METHODS.md}, and \texttt{info/STUDY\_SUMMARY.md}. Figure aggregates and source hashes accompany this report.
\end{thebibliography}\end{singlespace}
